\section{A minimal structural-biology vocabulary}
A protein is built from amino-acid residues connected along one or more chains. A residue contains several atoms; the C-alpha atom supplies one convenient reference point for its position along the backbone. Replacing a structure by its C-alpha coordinates produces a point cloud that retains the coarse arrangement of the resolved chain while discarding most atomic detail. An edge in a geometric graph on this cloud indicates a chosen spatial relationship. It need not be a covalent bond, and residues close in three-dimensional space need not be adjacent in sequence.

A deposited coordinate model is also not a movie. The crystallographic B-factor attached to an atom is an atomic displacement parameter, interpreted within the structure-determination model. It can be influenced by motion and disorder, but the coordinate file alone does not separate those contributions. This is why the study predicts a stated crystallographic observable rather than claiming to recover solution trajectories. The formula relating B-factors to isotropic displacement appears in the next section; the experimental outcome below is further standardized within each protein.

The three views in Figure~\ref{fig:protein} can now be read without knowing any sheaf theory. The first is the measured relative profile mapped to the structure. The second is a prediction for a protein excluded from that model's training groups. Their common color scale permits a direct visual comparison. The third uses a different scale because it shows only the signed contribution of one feature family to the prediction. A warm region in that panel says that this group raises the fitted model's output relative to its background expectation. It does not say that a sheaf feature causes motion at that residue.

